\section{Materials and methods}

\subsection{Backbone and masked gene expression}

The input is a processed gene expression matrix $X \in \mathbb{R}^{C \times G}$, where
$X_{ij}$ is the expression of gene $j$ in cell $i$, $C$ is the number of cells, and $G$
the number of genes. The backbone follows scMAE~\citep{scmae2024}: an encoder maps a
corrupted input to a latent embedding, a mask predictor infers which entries were
corrupted, and a decoder reconstructs the original expression.

Corruption is generated by a Bernoulli mask,
\begin{equation}
M_{ij} \sim \mathrm{Bernoulli}(p_j),
\end{equation}
where $p_j$ controls the fraction of gene $j$ that is corrupted. In contrast to the
per-gene value shuffling of the original scMAE, we corrupt by zeroing,
\begin{equation}
X^{M}_{ij} = X_{ij}\,(1 - M_{ij}).
\end{equation}
This choice is deliberate. On fine-grained datasets, global value shuffling is
destructive: on Macosko it yields an ARI of $0.30$ against $0.70$ for zero-masking under
an otherwise identical pipeline. Zero-masking is also the augmentation reported to be
most effective for single-cell self-supervised learning~\citep{scsslbench2025}. A masking
probability in the range $0.2$--$0.4$ is stable, consistent with the original
sensitivity analysis.

\subsection{Mask prediction and weighted reconstruction}

The encoder produces the embedding $E$. The mask predictor, a linear layer, estimates the
mask $\widetilde{M}$ from $E$ and is trained with cross-entropy,
\begin{equation}
L_m = -\sum_{ij} M_{ij}\log \widetilde{M}_{ij}.
\end{equation}
The decoder takes the concatenation of $E$ and $\widetilde{M}$ and reconstructs $X$ with a
weighted mean squared error,
\begin{equation}
L_r = \frac{1}{N}\sum_{ij}\Omega_{ij}\,\bigl(X_{ij}-\widetilde{X}_{ij}\bigr)^2,
\end{equation}
\begin{equation}
\Omega_{ij} = M_{ij}\,\lambda + (1-M_{ij})\,(1-\lambda),
\end{equation}
so that corrupted entries receive weight $\lambda$ and clean entries $1-\lambda$. The
backbone loss combines the two terms,
\begin{equation}
L_{\mathrm{scMAE}} = (1-\gamma)\,L_r + \gamma\,L_m.
\end{equation}
The mask-to-decoder path is retained because, in a controlled ablation, removing it does
not act as a shortcut to be discarded but as a stabilizer: without it the mean ARI on
Macosko drops and the seed variance grows.

\subsection{Coupled deep embedded clustering}

To couple representation and clustering---the property whose absence is identified as a
weakness of most methods~\citep{scclubench2025}---we place trainable cluster centers
$\mu_k$ in the latent space and compute a Student's $t$ soft assignment,
\begin{equation}
q_{ik} = \frac{\bigl(1+\lVert z_i-\mu_k\rVert^2\bigr)^{-1}}
{\sum_{k'}\bigl(1+\lVert z_i-\mu_{k'}\rVert^2\bigr)^{-1}},
\end{equation}
where $z_i$ is the latent code of cell $i$. Following the deep embedded clustering
objective~\citep{dec2016}, a sharpened target $p_{ik}$ is formed from $q$, and the
clustering loss is the divergence
\begin{equation}
L_{\mathrm{KL}} = \sum_{i}\sum_{k} p_{ik}\log\frac{p_{ik}}{q_{ik}},
\end{equation}
applied only after a warmup period and only to cells whose assignment confidence exceeds
a threshold. The warmup prevents early, unreliable centers from distorting the embedding;
the confidence gate restricts sharpening to well-formed assignments. These settings
reproduce a known configuration and are held fixed throughout.

\subsection{Diagnosis and the per-dimension variance floor}

The sharpening in Eq.~(8) compacts clusters but, as a side effect, suppresses the variance
of individual latent coordinates. We quantify this with the effective dimensionality,
defined as the participation ratio of the per-dimension variance spectrum,
\begin{equation}
d_{\mathrm{eff}} = \frac{\bigl(\sum_d \sigma_d^2\bigr)^2}{\sum_d \sigma_d^4},
\end{equation}
where $\sigma_d^2$ is the variance of latent coordinate $d$ across cells. On Macosko,
$d_{\mathrm{eff}}$ falls from $114$ (of $128$ dimensions) under the backbone to $60$ once
DEC sharpening is active. To prevent this collapse we add a per-dimension variance floor,
\begin{equation}
L_{\mathrm{floor}} = \frac{1}{D}\sum_{d=1}^{D}\max\!\bigl(0,\;\tau - \sigma_d\bigr),
\end{equation}
with $\tau = 1$ and $\sigma_d$ the batch standard deviation of coordinate $d$. This term
is exactly the variance term of VICReg~\citep{vicreg2022}, a mechanism established as
effective in single-cell self-supervised learning~\citep{scsslbench2025}. The total
objective is
\begin{equation}
L = L_{\mathrm{scMAE}} + \lambda_{\mathrm{clu}}\,L_{\mathrm{KL}}
    + \lambda_{\mathrm{floor}}\,L_{\mathrm{floor}}.
\end{equation}

Three design decisions warrant explanation. First, the floor acts on each latent
coordinate independently because the failure is per-dimension: it restores
$d_{\mathrm{eff}}$ from $60$ back to $114$, whereas decorrelation or uniformity objectives,
which do not target per-coordinate variance, do not. Second, the floor is applied to the
entire clustering latent rather than to a reserved subspace, because its mechanism is to
maximize the number of active clustering dimensions; splitting the latent removes that
capacity and collapses performance to an ARI near $0.31$ in direct tests. Third, we add no
structure derived from the expression matrix, such as gene programs, neighbor graphs, or
marker sets. In an ablation, using a real non-negative matrix factorization program as an
auxiliary target yields an ARI of $0.458$, lower than a meaningless row-permuted target
($0.632$) and both below the baseline ($0.779$); expression-derived semantics pull the
embedding away from the cell-type axes. The design is therefore intentionally minimal:
backbone, coupled clustering, and a variance floor.

\subsection{Datasets, comparison methods, and evaluation}

We evaluate on 23 datasets spanning multiple tissues, organisms, and sequencing platforms,
including several with rare cell types (Macosko, Shekhar, Tosches, Baron, Quake Spleen).
We compare against representative methods from all families: the masked autoencoder scMAE
and its neighbor-mixing variant; the deep clustering methods DEC, scDeepCluster, scDCC,
DESC, scNAME, and scziDesk; the generative model scVI~\citep{scvi2018}; the graph method
AttentionAE-sc; the traditional methods Leiden, Louvain, and SC3; the foundation model
scGPT; and a plain PCA+$k$-means baseline with the number of clusters set to the true
value.

Evaluation uses four metrics rather than the customary three. In addition to the Adjusted
Rand Index (ARI) and Normalized Mutual Information (NMI), which reward correct assignment
of abundant cells, we report macro-averaged F1 and rare-cluster recall, defined as the
recall averaged over classes that account for less than $2\%$ of cells. ARI and accuracy
are dominated by large classes and can conceal failures on rare populations; the two
additional metrics make rare-cell performance explicit. The known-$k$ setting fixes the
number of clusters to the ground-truth value and does not constitute a separate algorithm.
Each dataset--method pair is run with three random seeds, and results are reported as the
mean.
